\ChapterDivider
  {Electrical Machines}
  {Magnetic Circuits}
    {\begin{enumerate}[leftmargin=*,itemsep=2pt,topsep=2pt]
    \item Magnetically Coupled Circuits, Dot Convention
    \item Energy in Coupled circuits
    \item Electric and Magnetic Materials
  \end{enumerate}}

  \MCQ{1}{2025}{1}{A, D}{2}{
Consider two coupled circuits, having self-inductances $L_1$ and $L_2$, that carry non-zero currents $I_1$ and $I_2$, respectively. The mutual inductance between the circuits is $M$ with unity coupling coefficient. The stored magnetic energy of the coupled circuits is minimum at which of the following value(s) of $\dfrac{I_1}{I_2}$?

\InlineOptionsOneLine
{$-\dfrac{M}{L_1}$}
{$-\dfrac{M}{L_2}$}
{$-\dfrac{L_1}{M}$}
{$-\dfrac{L_2}{M}$}
}
\MCQ{1}{2021}{1}{1}{A}{The input impedance, $Z_{in}(s)$, for the network shown is
\QuestionFigure{img/ch1_Magnetic_circuits/1_21_1.png}
\InlineOptions
{$\dfrac{23s^2 + 46s + 20}{4s + 5}$}
{$6s + 4$}
{$7s + 4$}
{$\dfrac{25s^2 + 46s + 20}{4s + 5}$}
}
\NAT{1}{2019}{1}{2}{0.83 to 0.84}{The magnetic circuit shown below has a uniform cross-sectional area and an air gap of $0.2 \ \text{cm}$. The mean path length of the core is $40 \ \text{cm}$. Assume that leakage and fringing fluxes are negligible. When the core relative permeability is assumed to be infinite, the magnetic flux density computed in the air gap is $1 \ \text{tesla}$. With the same Ampere-turns, if the core relative permeability is assumed to be $1000$ (linear), the flux density in tesla (rounded off to three decimal places) calculated in the air gap is \underline{\hspace{2cm}}.
}
\QuestionFigure{img/ch1_Magnetic_circuits/1_19_1.png}

\MCQ{1}{2018}{1}{2}{D}{A transformer with a toroidal core of permeability $\mu$ is shown in the figure. Assuming a uniform flux density across the circular core cross-section of radius $r \ll R$, and neglecting any leakage flux, the best estimate for the mean radius $R$ is

\QuestionFigure{img/ch1_Magnetic_circuits/1_18_1.png}

\InlineOptions
{$\dfrac{\mu V r^2 N_p^2 \omega}{I}$}
{$\dfrac{\mu I r^2 N_p N_s \omega}{V}$}
{$\dfrac{\mu V r^2 N_p^2 \omega}{2I}$}
{$\dfrac{\mu V r^2 N_p^2 \omega}{2V}$}
}
\NAT{1}{2016}{1}{2}{186 to 190}{
The flux linkage $(\lambda)$ and current $(i)$ relation for an electromagnetic system is

$\lambda=\dfrac{\sqrt{i}}{g}$.

When $i=2\,\text{A}$ and $g$ (air-gap length) $=10\,\text{cm}$, the magnitude of mechanical force on the moving part, in N, is \rule{2cm}{0.15mm}.
}
\NAT{1}{2015}{1}{2}{0}{
Two identical coils each having inductance $L$ are placed together on the same core. If an overall inductance of $\alpha L$ is obtained by interconnecting these two coils, the minimum value of $\alpha$ is \rule{2cm}{0.15mm}.
}
\CommonData{
\textbf{Statement for Linked  Questions (Q1.07.1 and Q1.07.2):}
An inductor designed with $400$ turns coil wound on an iron core of $16\,\text{cm}^2$ cross sectional area and with a cut of an air gap length of $1\,\text{mm}$. The coil is connected to a $230\,\text{V}$, $50\,\text{Hz}$ ac supply. Neglect coil resistance, core loss, iron reluctance and leakage inductance. $(\mu_0 = 4\pi \times 10^{-7}\,\text{H/m})$
}

\MCQ{1}{2007}{1}{2}{D}{
The current in the inductor is

\InlineOptionsOneLine{18.08\,A}{9.04\,A}{4.56\,A}{2.28\,A}
}

\MCQ{1}{2007}{2}{2}{A}{
The average force on the core to reduce the air gap will be

\InlineOptionsOneLine{832.29\,N}{1666.22\,N}{3332.47\,N}{6664.84\,N}
}
\MCQ{1}{2004}{1}{1}{A}{
For a linear electromagnetic circuit, the following statement is true.

\begin{choices}
\Option{Field energy is equal to the co-energy}
\Option{Field energy is greater than the co-energy}
\Option{Field energy is lesser than the co-energy}
\Option{Co-energy is zero}
\end{choices}
}
\NAT{1}{2001}{1}{5}{530.52, 10}{
Determine the resonance frequency and the Q-factor of the circuit shown in figure. 

Data: $R = 10\,\Omega$, $C = 3\,\mu\text{F}$, $L_1 = 40\,\text{mH}$, $L_2 = 10\,\text{mH}$ and $M = 10\,\text{mH}$.
\QuestionFigure{img/ch1_Magnetic_circuits/1_01_1.png}
}

\MCQ{1}{2001}{2}{2}{C}{The hysteresis loop of a magnetic material has an area of $5\text{ cm}^{2}$, with the scales given as $1\text{ cm} = 2\text{ AT}$ and $1\text{ cm} = 50\text{ mWb}$. At $50\text{ Hz}$, the total hysteresis loss is
\InlineOptionsOneLine{$15\text{ W}$}{$20\text{ W}$}{$25\text{ W}$}{$50\text{ W}$}}

\NAT{1}{1997}{1}{5}{0.0628}{}{The given figure shows a magnetic circuit formed by an ideal core material. Determine the magnetic flux density in the air gap.
\NATBox}
\QuestionFigure{img/ch1_Magnetic_circuits/1_97_3.png}

\NAT{1}{1997}{2}{2}{0.11}{Two identical coils of negligible resistance, when connected in series across a $200\,\text{V}$, $50\,\text{Hz}$ source, draw a current of $10\,\text{A}$. When the terminals of one of the coils are reversed, the current drawn is $8\,\text{A}$. The coefficient of coupling between the two coils is \underline{\hspace{2cm}}.
}

\NAT{1}{1995}{1}{5}{303.3, 0414}%
{Figure below, shows the details of an electromagnet with an operating coil rated for $230\,\text{V}$, $50\,\text{Hz}$. The coil has $1500$ turns. Find the average force of attraction on the armature of the electromagnet and the current drawn by the coil. Omit winding resistance and core losses and assume infinite permeability for the core.

\QuestionFigure[0.9\columnwidth]{img/ch1_Magnetic_circuits/1_95_1.png}
}
\MCQ{1}{1992}{1}{1}{B}%
{The developed electromagnetic force and/or torque in electro-mechanical energy conversion systems act in a direction that tends

\begin{choices}
\Option{to increase the stored energy at constant $mmf$}
\Option{to decrease the stored energy at constant flux}
\Option{to decrease the co-energy at constant $mmf$}
\Option{to decrease the stored energy at constant $mmf$}
\end{choices}
}
\MCQ{1}{1991}{1}{2}{
$100,\; 1.44 \times 10^7,\; 0.69,\; 0.345$
}{A toroidal iron ring has a uniform cross-sectional area of 50 mm² and a mean magnetic path length of 100 mm. The ring has an air-gap of 1 mm, and it is excited with a dc current of 1 A through a coil of 100 turns wound uniformly along its length. The iron may be assumed to be a perfect magnetic material. The effect of fringing at the gap may be assumed to increase the effective area of magnetic flux at the gap by 10\%. Evaluate
\begin{enumerate}
\item[(a)] The exciting mmf of the coil
\item[(b)] The effective reluctance of the magnetic circuit.
\item[(c)] The inductance of the coil
\item[(d)] The energy stored in the magnetic field under the above excitation.
\end{enumerate}
}
\MCQ{1}{1991}{2}{1}{B}%
{The energy stored in the magnetic field of solenoid $30\,\text{cm}$ long and $3\,\text{cm}$ diameter wound with $1000$ turns of wire carrying a current of $10\,\text{A}$ is

\InlineOptions
{$0.015\,\text{joule}$}
{$0.15\,\text{joule}$}
{$0.5\,\text{joule}$}
{$1.15\,\text{joule}$}
}
\EndChapter